\documentclass[a4paper,landscape,8pt]{extarticle}
\usepackage[landscape,margin=0.7cm,top=0.5cm,bottom=0.7cm]{geometry}
\usepackage[utf8]{inputenc}
\usepackage[T1]{fontenc}
\usepackage[ngerman]{babel}
\usepackage{helvet}
\renewcommand{\familydefault}{\sfdefault}
\usepackage{xcolor}
\usepackage{tikz}
\usetikzlibrary{positioning}
\usepackage{enumitem}
\usepackage{multicol}
\usepackage{array}
\usepackage{fancyhdr}
\usepackage{amssymb}

% Farben
\definecolor{mailblue}{HTML}{1565C0}
\definecolor{foldergreen}{HTML}{2E7D32}
\definecolor{attachorange}{HTML}{E65100}
\definecolor{spamred}{HTML}{C62828}
\definecolor{lightgray}{HTML}{F5F5F5}
\definecolor{bordergray}{HTML}{BBBBBB}

\pagestyle{fancy}
\fancyhf{}
\renewcommand{\headrulewidth}{0pt}
\fancyfoot[C]{\footnotesize Seite \thepage{} von 3}

\setlength{\parindent}{0pt}
\setlength{\parskip}{2pt}

% Kompakte Listen
\setlist{nosep, leftmargin=1em, itemsep=0pt, topsep=1pt}

\begin{document}

% ============================================================================
% SEITE 1: INFORMATIONSBLATT
% ============================================================================

\noindent{\Large\bfseries E-Mail verstehen: Dein digitaler Briefkasten} \hfill
{\footnotesize\textbf{Lernziele:} Ordner kennen | Mails finden | Anhänge öffnen}

\vspace{1mm}
\hrule
\vspace{1mm}

% Drei-Spalten-Layout
\begin{minipage}[t]{0.32\linewidth}
\begin{center}
\colorbox{mailblue!10}{%
\begin{minipage}{0.95\linewidth}
\vspace{2mm}
\begin{center}
{\Large\bfseries\textcolor{mailblue}{1. Die E-Mail-Adresse}}\\[1pt]
{\footnotesize\itshape Deine digitale Postadresse}
\end{center}
\vspace{2mm}
\end{minipage}%
}
\end{center}

\vspace{1mm}

\textbf{So sieht sie aus:}\\
\texttt{vorname.nachname@anbieter.de}

\textbf{Die drei Teile:}
\begin{itemize}
\item \textbf{Vor dem @:} Dein Name (frei wählbar)
\item \textbf{Das @:} ``bei'' -- trennt Name von Anbieter
\item \textbf{Nach dem @:} Der Anbieter (gmail.com, web.de, gmx.de...)
\end{itemize}

\textbf{Analogie:} Wie eine Postadresse -- damit Nachrichten bei dir ankommen.

\textbf{Wichtig:}
\begin{itemize}
\item Groß-/Kleinschreibung egal
\item Keine Leerzeichen
\item Genau so tippen wie vorgegeben
\end{itemize}

\end{minipage}%
\hfill%
\begin{minipage}[t]{0.32\linewidth}
\begin{center}
\colorbox{foldergreen!10}{%
\begin{minipage}{0.95\linewidth}
\vspace{2mm}
\begin{center}
{\Large\bfseries\textcolor{foldergreen}{2. Die wichtigen Ordner}}\\[1pt]
{\footnotesize\itshape Wo sind meine Mails?}
\end{center}
\vspace{2mm}
\end{minipage}%
}
\end{center}

\vspace{1mm}

\textbf{Posteingang} (Inbox)
\begin{itemize}
\item Hier landen \textbf{alle neuen} Mails
\item Wie dein Briefkasten zu Hause
\end{itemize}

\textbf{Gesendet}
\begin{itemize}
\item Kopien deiner \textbf{verschickten} Mails
\item Wie Durchschläge deiner Briefe
\end{itemize}

\textbf{Entwürfe}
\begin{itemize}
\item Angefangene, \textbf{noch nicht} gesendete Mails
\end{itemize}

\textbf{Papierkorb}
\begin{itemize}
\item Gelöschte Mails (noch \textbf{30 Tage} rettbar)
\item Danach endgültig weg
\end{itemize}

\textbf{Spam/Werbung}
\begin{itemize}
\item Automatisch \textbf{aussortierte} verdächtige Mails
\item Manchmal landen hier auch \textbf{echte} Mails!
\end{itemize}

\end{minipage}%
\hfill%
\begin{minipage}[t]{0.32\linewidth}
\begin{center}
\colorbox{attachorange!10}{%
\begin{minipage}{0.95\linewidth}
\vspace{2mm}
\begin{center}
{\Large\bfseries\textcolor{attachorange}{3. Anhänge öffnen}}\\[1pt]
{\footnotesize\itshape Dateien in der Mail}
\end{center}
\vspace{2mm}
\end{minipage}%
}
\end{center}

\vspace{1mm}

\textbf{Was ist ein Anhang?}\\
Dateien, die jemand der Mail \textbf{beigelegt} hat -- wie ein Foto im Briefumschlag.

\textbf{So erkennst du Anhänge:}
\begin{itemize}
\item \textbf{Büroklammer-Symbol} in der Mail-Liste
\item Dateien \textbf{unter dem Text} der Mail
\item Vorschau-Bilder bei Fotos
\end{itemize}

\textbf{So öffnest du sie:}
\begin{enumerate}
\item Mail antippen/öffnen
\item Nach unten scrollen
\item Auf die Datei tippen
\item Vorschau erscheint
\end{enumerate}

\textbf{So speicherst du:}
\begin{itemize}
\item Teilen-Symbol antippen
\item ``In Dateien sichern'' wählen
\end{itemize}

\vspace{2mm}
\fcolorbox{spamred}{spamred!5}{%
\begin{minipage}{0.92\linewidth}
\footnotesize
\textbf{Vorsicht:} Anhänge von Unbekannten NICHT öffnen! (Viren-Gefahr)
\end{minipage}%
}

\end{minipage}

\vspace{1mm}

% Zusammenfassung
\begin{center}
\fcolorbox{bordergray}{lightgray}{%
\begin{minipage}{0.98\linewidth}
\centering
\vspace{1.5mm}
\textbf{Zusammenfassung:}
\textcolor{mailblue}{\textbf{E-Mail-Adresse}} = deine digitale Postadresse ~$|$~
\textcolor{foldergreen}{\textbf{Ordner}} = Posteingang, Gesendet, Papierkorb, Spam ~$|$~
\textcolor{attachorange}{\textbf{Anhänge}} = beigelegte Dateien (Büroklammer)
\vspace{1.5mm}
\end{minipage}%
}
\end{center}

\newpage

% ============================================================================
% SEITE 2: ÜBUNGEN
% ============================================================================

\begin{center}
{\LARGE\bfseries Übungen: Teste dein Wissen}\\[3pt]
{\normalsize Schau dir Seite 1 an und beantworte die Fragen}
\end{center}

\vspace{3mm}
\hrule
\vspace{4mm}

\begin{multicols}{2}

\textbf{\large Teil A: E-Mail-Adressen verstehen}

Welcher Teil ist was? Verbinde mit Linien.

\vspace{2mm}

\begin{center}
\texttt{maria.schmidt@gmx.de}
\end{center}

\vspace{2mm}

\renewcommand{\arraystretch}{1.8}
\begin{tabular}{@{}p{4cm}p{4cm}@{}}
\texttt{maria.schmidt} & $\circ$ ~~~~ $\circ$ Der Anbieter \\
\texttt{@} & $\circ$ ~~~~ $\circ$ Dein gewählter Name \\
\texttt{gmx.de} & $\circ$ ~~~~ $\circ$ ``bei'' (trennt Name und Anbieter) \\
\end{tabular}
\renewcommand{\arraystretch}{1}

\vspace{6mm}

\textbf{\large Teil B: Ordner zuordnen}

In welchen Ordner gehört das? Schreibe: \textit{Posteingang}, \textit{Gesendet}, \textit{Papierkorb} oder \textit{Spam}

\vspace{2mm}

\renewcommand{\arraystretch}{1.8}
\begin{tabular}{@{}p{8cm}p{4cm}@{}}
1. Eine Mail von deiner Tochter kommt an: & \dotfill \\
2. Eine Mail, die du geschrieben hast: & \dotfill \\
3. Eine Mail, die du gelöscht hast: & \dotfill \\
4. Eine verdächtige Werbe-Mail: & \dotfill \\
5. Du suchst eine alte Mail, die du geschickt hast: & \dotfill \\
\end{tabular}
\renewcommand{\arraystretch}{1}

\vspace{6mm}

\textbf{\large Teil C: Richtig oder Falsch?}

Kreuze an: R = richtig, F = falsch

\vspace{2mm}

\renewcommand{\arraystretch}{1.6}
\begin{tabular}{@{}p{9cm}cc@{}}
& R & F \\
1. Im Papierkorb sind Mails sofort weg. & $\square$ & $\square$ \\
2. Anhänge erkennt man am Büroklammer-Symbol. & $\square$ & $\square$ \\
3. Der Spam-Ordner fängt nur schlechte Mails. & $\square$ & $\square$ \\
4. Groß- und Kleinschreibung ist bei E-Mail-Adressen wichtig. & $\square$ & $\square$ \\
5. Eine Mail kann mehrere Anhänge haben. & $\square$ & $\square$ \\
\end{tabular}
\renewcommand{\arraystretch}{1}

\columnbreak

\textbf{\large Teil D: Situationen}

Was tust du? Schreibe die Antwort in Stichworten.

\vspace{3mm}

\textbf{Situation 1:}\\
Du erwartest eine wichtige Mail, aber sie ist nicht im Posteingang.

\textit{Was machst du?}
\vspace{6mm}

\rule{\linewidth}{0.4pt}
\vspace{5mm}

\rule{\linewidth}{0.4pt}

\vspace{5mm}

\textbf{Situation 2:}\\
Jemand schickt dir ein Foto per Mail. Wie öffnest du es?

\textit{Deine Schritte:}
\vspace{6mm}

\rule{\linewidth}{0.4pt}
\vspace{5mm}

\rule{\linewidth}{0.4pt}

\vspace{5mm}

\textbf{Situation 3:}\\
Du hast versehentlich eine wichtige Mail gelöscht.

\textit{Was kannst du tun?}
\vspace{6mm}

\rule{\linewidth}{0.4pt}
\vspace{5mm}

\rule{\linewidth}{0.4pt}

\vspace{5mm}

\textbf{Situation 4:}\\
Du bekommst eine Mail mit Anhang von einer unbekannten Adresse.

\textit{Was solltest du tun?}
\vspace{6mm}

\rule{\linewidth}{0.4pt}

\end{multicols}

\newpage

% ============================================================================
% SEITE 3: CHECKLISTE
% ============================================================================

\begin{center}
{\LARGE\bfseries Checkliste: Das habe ich verstanden}\\[3pt]
{\normalsize Geh die Liste durch und hake ab, was du sicher weißt}
\end{center}

\vspace{3mm}
\hrule
\vspace{6mm}

% Drei Boxen nebeneinander
\begin{minipage}[t]{0.31\linewidth}
\begin{center}
\fcolorbox{mailblue}{mailblue!8}{%
\begin{minipage}{0.92\linewidth}
\vspace{3mm}
\begin{center}
{\large\bfseries\textcolor{mailblue}{E-Mail-Grundlagen}}
\end{center}
\vspace{2mm}
\begin{itemize}[label=$\square$, itemsep=5pt]
\item Ich verstehe den \textbf{Aufbau} einer E-Mail-Adresse
\item Ich weiß, was das \textbf{@-Zeichen} bedeutet
\item Ich kann die \textbf{Mail-App} öffnen
\item Ich finde den \textbf{Posteingang}
\end{itemize}
\vspace{3mm}
\end{minipage}%
}
\end{center}
\end{minipage}%
\hfill%
\begin{minipage}[t]{0.31\linewidth}
\begin{center}
\fcolorbox{foldergreen}{foldergreen!8}{%
\begin{minipage}{0.92\linewidth}
\vspace{3mm}
\begin{center}
{\large\bfseries\textcolor{foldergreen}{Ordner verstehen}}
\end{center}
\vspace{2mm}
\begin{itemize}[label=$\square$, itemsep=5pt]
\item Ich kenne den \textbf{Posteingang}
\item Ich finde \textbf{gesendete} Mails
\item Ich weiß, wo der \textbf{Papierkorb} ist
\item Ich prüfe manchmal den \textbf{Spam}-Ordner
\end{itemize}
\vspace{3mm}
\end{minipage}%
}
\end{center}
\end{minipage}%
\hfill%
\begin{minipage}[t]{0.31\linewidth}
\begin{center}
\fcolorbox{attachorange}{attachorange!8}{%
\begin{minipage}{0.92\linewidth}
\vspace{3mm}
\begin{center}
{\large\bfseries\textcolor{attachorange}{Anhänge}}
\end{center}
\vspace{2mm}
\begin{itemize}[label=$\square$, itemsep=5pt]
\item Ich erkenne Mails \textbf{mit Anhang}
\item Ich kann Anhänge \textbf{öffnen}
\item Ich kann Anhänge \textbf{speichern}
\item Ich öffne \textbf{keine} Anhänge von Unbekannten
\end{itemize}
\vspace{3mm}
\end{minipage}%
}
\end{center}
\end{minipage}

\vspace{8mm}

% Gesamtverständnis
\begin{center}
\fcolorbox{black}{lightgray}{%
\begin{minipage}{0.95\linewidth}
\vspace{4mm}
\begin{center}
{\large\bfseries Gesamtverständnis}
\end{center}
\vspace{3mm}
\begin{multicols}{2}
\begin{itemize}[label=$\square$, itemsep=6pt]
\item Ich kann mich in der \textbf{Mail-App orientieren}
\item Ich finde Mails in den \textbf{verschiedenen Ordnern}
\item Ich kann \textbf{Anhänge} öffnen und speichern
\item Ich bin \textbf{vorsichtig} bei unbekannten Absendern
\item Ich fühle mich \textbf{sicherer} im Umgang mit E-Mails
\end{itemize}
\end{multicols}
\vspace{3mm}
\end{minipage}%
}
\end{center}

\vspace{8mm}

% Notfall-Notizen
\begin{center}
\fcolorbox{bordergray}{white}{%
\begin{minipage}{0.95\linewidth}
\vspace{4mm}
\begin{center}
{\large\bfseries Meine E-Mail-Notizen}\\[3pt]
{\footnotesize (Fülle das hier aus und bewahre es sicher auf!)}
\end{center}
\vspace{4mm}
\renewcommand{\arraystretch}{2}
\begin{tabular}{@{}p{0.45\linewidth}p{0.45\linewidth}@{}}
\textbf{Meine E-Mail-Adresse:} & \textbf{Mein E-Mail-Anbieter:} \\
\dotfill & \dotfill \\[3mm]
\textbf{Wer hilft mir bei Problemen?} & \\
\dotfill & \\
\end{tabular}
\renewcommand{\arraystretch}{1}
\vspace{4mm}
\end{minipage}%
}
\end{center}

\vfill

\begin{center}
\footnotesize\textit{Stand: \today{} -- Öffne niemals Anhänge von unbekannten Absendern!}
\end{center}

\end{document}
